\begin{textsample}{POS Dim 3 – rj – Score 18.00 – rj/rj\_inf\_curriculo\_1.txt}  \label{ex:f3_pos_017}
A Proposta Curricular da Educação \textbf{Infantil} Planejar é preciso! \textbf{Criança}, centro do planejamento curricular, é sujeito histórico e de direitos que, nas interações, relações e práticas cotidianas que vivencia, constrói sua identidade pessoal e coletiva, \textbf{brinca}, imagina, fantasia, deseja, aprende, observa, experimenta, narra, questiona e constrói sentidos sobre a natureza e a sociedade, produzindo cultura ( BRASIL, 2009, Art. 4º ).
Planejar é um ato de \textbf{cuidado} com a prática pedagógica destinada aos \textbf{pequenos}.
O planejamento do professor de Educação \textbf{Infantil} é uma proposta que reorganiza e norteia todas as ações educativas que envolvem as \textbf{crianças}, tanto no âmbito pedagógico propriamente dito, quanto nas inter-relações de todos os envolvidos nesse processo, dando voz às \textbf{crianças} e \textbf{acolhendo} a forma como elas significam o mundo e a si mesmas.
É necessário explicitar os campos de experiências de modo que o cotidiano apareça sem rigidez e sempre transparente e flexível.
Currículo e práticas pedagógicas Um dos grandes desafios da Educação \textbf{Infantil} é a busca pelos conhecimentos sociais relevantes e significativos, trabalhando -se não só com a imaginação, mas também com observações, comparações e com autocrítica.
Em suma, para se desenvolver, aprender e construir conhecimentos, a \textbf{criança} precisa se expressar por meio de diferentes linguagens, agir, perguntar, ler o mundo, olhar imagens, criar relações, \textbf{testar} hipóteses e refletir sobre o que faz ou aprende, de modo a reestruturar o pensamento permanentemente.
As práticas pedagógicas que compõem a proposta curricular da Educação \textbf{Infantil} tem como eixos norteadores as interações e a \textbf{brincadeira}, garantindo à \textbf{criança} os direitos de desenvolvimento e aprendizagem.
Cabe salientar as Diretrizes Curriculares Nacionais da Educação \textbf{Infantil} em seu artigo 9º : Art. 9º : As práticas pedagógicas que compõem a proposta curricular da Educação \textbf{Infantil} devem ter como eixos norteadores as interações e a \textbf{brincadeira}, garantindo experiências que : I - promovam o conhecimento de si e do mundo por meio da ampliação de experiências \textbf{sensoriais}, expressivas, corporais que possibilitem \textbf{movimentação} ampla, expressão da individualidade e respeito pelos ritmos e \textbf{desejos} da \textbf{criança} ; II - favoreçam a \textbf{imersão} das \textbf{crianças} nas diferentes linguagens e o progressivo domínio por elas de vários gêneros e formas de expressão : gestual, verbal, \textbf{plástica}, \textbf{dramática} e musical ; III - possibilitem às \textbf{crianças} experiências de narrativas, de apreciação e interação com a linguagem oral e escrita, e convívio com diferentes suportes e gêneros textuais orais e escritos ; IV - recriem, em contextos significativos para as \textbf{crianças}, relações quantitativas, medidas, formas e orientações espaço-temporais ; V - ampliem a confiança e a participação das \textbf{crianças} nas atividades individuais e coletivas ; VI - possibilitem situações de aprendizagem mediadas para a elaboração da autonomia das \textbf{crianças} nas ações de \textbf{cuidado} pessoal, auto-organização, saúde e bem-estar ; VII - possibilitem vivências éticas e estéticas com outras \textbf{crianças} e grupos culturais, que alarguem seus padrões de referência e de identidades no diálogo e reconhecimento da diversidade ; VIII - incentivem a \textbf{curiosidade}, a exploração, o encantamento, o questionamento, a indagação e o conhecimento das \textbf{crianças} em relação ao mundo físico e social, ao tempo e à natureza ; IX - promovam o relacionamento e a interação das \textbf{crianças} com diversificadas manifestações de música, artes \textbf{plásticas} e gráficas, cinema, fotografia, dança, teatro, poesia e literatura ; X - promovam a interação, o \textbf{cuidado}, a preservação e o conhecimento da biodiversidade e da sustentabilidade da vida na Terra, assim como o não desperdício dos recursos naturais ; XI - propiciem a interação e o conhecimento pelas \textbf{crianças} das manifestações e tradições culturais brasileiras ; XII - possibilitem a utilização de gravadores, projetores, computadores, máquinas \textbf{fotográficas}, e outros recursos tecnológicos e midiáticos.
Parágrafo único - As \textbf{creches} e \textbf{pré-escolas}, na elaboração da proposta curricular, de acordo com suas características, identidade institucional, escolhas coletivas e particularidades pedagógicas, estabelecerão modos de integração dessas experiências.

% matched lemmas: acolher, brincadeira, brincar, creche, criança, cuidado, curiosidade, desejo, dramático, fotográfico, imersão, infantil, movimentação, pequeno, plástico, pré-escola, sensorial, testar
\end{textsample}
